%! TEX root = definisi_tapak_temu1.tex
% the main TeX file which is intended to compile, :VimtexReload after adjustment
% :h vimtex-tex-root
\documentclass[../main.tex]{subfiles}
% \includeonlyframes{current}
% \setbeameroption{hide notes} % Only slides
% \setbeameroption{show only notes} % Only notes
% \setbeameroption{show notes on second screen=right} % Both

\begin{document}

\begin{frame}[mytitle=standard,t,mycolor=digiPH_gray,light]
	\begin{columns}
		\begin{column}{0.5\textwidth}
			Kelompok 1
			\begin{enumerate}
				\item \textbf{Muh. Aril}
				\item Andi Rifat Yusuf
				\item Mutiara Salim
				\item Firman
				\item Wahyuni
			\end{enumerate}
			Kelompok 3
			\begin{itemize}
				\item \textbf{Muthmainnah Humairo Idris}
				\item Ikmal Jalil
				\item Muhammad Firjatullah
				\item Muhammad Taufiq Ismail
				\item Siti Radiah Syah
			\end{itemize}
			Kelompok 5
			\begin{itemize}
				\item \textbf{Muhammad Fathir Rahmat}
				\item Keisha Tirta Adzani
				\item Afrilia Zhaifathun Qarin
				\item Sekar Ayu Suci Ramadhani
				\item Ch Syakila Nazza Rivai
			\end{itemize}



		\end{column}
		\begin{column}{0.5\textwidth}
			Kelompok 2
			\begin{itemize}
				\item \textbf{Putri Adiba Shofiyah}
				\item Mery Sapnah
				\item Alwitinro Amiruddin
				\item Indah Rahma Utami R.
				\item Suci Aulia N.
			\end{itemize}
			Kelompok 4
			\begin{itemize}
				\item \textbf{Andini Niramrrynur}
				\item Andi Ifan Trifajri
				\item Mahardika
				\item Azizah Putri Salsabila
				\item Clarisa Pasang
			\end{itemize}
			Kelompok 6
			\begin{itemize}
				\item \textbf{Ardian Gahral Saputra}
				\item Nurul Ramadhani Jalil
				\item Muh. Nur Hasby
				\item Indo Ungke Nurlia
				\item Nur Fadhilah Awal
			\end{itemize}


		\end{column}
	\end{columns}
\end{frame}


\begin{frame}[mytitle=standard,t,mycolor=digiPH_ocean,dark]
	\begin{columns}
		\begin{column}{0.5\textwidth}
			Kelompok 7
			\begin{itemize}
				\item \textbf{Qays Ahmad Deedat}
				\item Muhammad Rafli Irsan
				\item Muhammad Ficky Ramadhan
				\item Dwi Ayu Yulianti
				\item Riza Ramdani
			\end{itemize}
		\end{column}
		\begin{column}{0.5\textwidth}
			Kelompok 8
			\begin{itemize}
				\item \textbf{Yeheskiel Wilson Tarru}
				\item Natasya Pratiwi
				\item A. M. Zayyan Fahrezi
				\item Mohammad Syahrin
				\item Rizky Ananda Heriadi
				\item Emanuel Japakaimu Tokio
			\end{itemize}
		\end{column}
	\end{columns}


\end{frame}

\begin{frame}[mytitle=standard,t,mycolor=digiPH_uniblue,light]



\end{frame}



% ==========================================
% TOPIK 1: DEFINISI PERANCANGAN TAPAK
% ==========================================
\begin{frame}[mytitle=center, c, mycolor=digiPH_darkblue, dark]
	\frametitle{Definisi Perancangan Tapak}
	\framesubtitle{Perancangan Tapak (Site Planning)}

	Memahami tapak sebagai kesatuan ruang, fungsi, lingkungan, dan waktu.

	\vspace{0.3cm}
	\textbf{Subtopik utama:}
	\begin{itemize}
		\item Pengertian perancangan tapak
		\item Tujuan dan ruang lingkup
		\item Kaidah perancangan
		\item Unsur-unsur perencanaan tapak
		\item Faktor alam dan kondisi eksisting
		\item Tanah, topografi, angin, serta konsep cut and fill
	\end{itemize}

	\vspace{0.5cm}
	\textit{“Apakah bangunan dapat dirancang dengan baik tanpa memahami tempat di mana bangunan itu akan berdiri?”}
\end{frame}

% ==========================================
% TOPIK 2: CAPAIAN PEMBELAJARAN
% ==========================================
\begin{frame}[mytitle=standard, t, mycolor=digiPH_cream, light]
	\frametitle{Capaian Pembelajaran}

	Setelah mempelajari materi ini, pembelajar diharapkan mampu:
	\begin{itemize}
		\item Memahami pengertian dan ruang lingkup perancangan tapak.
		\item Mengidentifikasi potensi dan kendala suatu tapak.
		\item Menempatkan objek fisik dan kegiatan dalam satu kesatuan ruang dan waktu.
		\item Mengintegrasikan aspek arsitektur, teknik, arsitektur lanskap, dan perencanaan kota.
		\item Menganalisis orientasi, aksesibilitas, vegetasi, view, dan kondisi eksisting.
		\item Mengidentifikasi faktor alam yang memengaruhi perancangan tapak.
	\end{itemize}
\end{frame}

% ==========================================
% TOPIK 3: APA YANG DIMAKSUD DENGAN TAPAK?
% ==========================================
\begin{frame}[mytitle=standard, c, mycolor=digiPH_gray, light]
	\frametitle{Apa yang Dimaksud dengan Tapak?}

	Tapak merupakan lahan atau bagian ruang tempat berbagai kebutuhan bangunan, kegiatan, sirkulasi, ruang luar, dan elemen pendukung ditempatkan.

	\vspace{0.3cm}
	\textbf{Dalam konteks perancangan:} \\
	Tapak bukan sekadar ``tanah kosong''. Tapak memiliki:
	\begin{itemize}
		\item kondisi fisik dan kondisi alam,
		\item bangunan dan elemen eksisting,
		\item hubungan dengan lingkungan sekitar,
		\item potensi dan keterbatasan,
		\item serta karakter sosial dan aktivitas manusia.
	\end{itemize}

	\vspace{0.3cm}
	\begin{alertblock}{Inti pemahaman}
		Tapak adalah sistem ruang yang harus dibaca sebelum dirancang.
	\end{alertblock}
\end{frame}

% ==========================================
% TOPIK 4: DEFINISI PERANCANGAN TAPAK (LANJUTAN)
% ==========================================
\begin{frame}[mytitle=standard, t, mycolor=digiPH_white, light]
	\frametitle{Definisi Perancangan Tapak}

	\textbf{Rancangan Tapak (Site Plan):}
	\begin{quote}
		Gambaran atau peta rencana peletakan bangunan atau kavling beserta unsur penunjangnya dalam batas dan luas lahan tertentu.
	\end{quote}

	\textbf{Perencanaan Tapak:}
	\begin{quote}
		Pengolahan fisik tapak untuk meletakkan seluruh kebutuhan rancangan di dalam tapak.
	\end{quote}

	\vspace{0.3cm}
	Perencanaan dilakukan dengan memperhatikan kondisi tapak, kemungkinan dampak perubahan fisik, lingkungan alam, lingkungan fisik buatan, dan lingkungan sosial.
\end{frame}

% ==========================================
% TOPIK 5: PROSES INTEGRATIF
% ==========================================
\begin{frame}[mytitle=center, c, mycolor=digiPH_ocean, dark]
	\frametitle{Perancangan Tapak Sebagai Proses Integratif}

	Perancangan tapak tidak berdiri sendiri. Penggunaan lahan berkaitan dengan beberapa bidang:

	\begin{description}
		\item[ARSITEKTUR] Bangunan dan kavling.
		\item[TEKNIK] Jalan, drainase, energi, limbah, dan utilitas.
		\item[ARSITEKTUR LANSKAP] Ruang terbuka hijau dan nonhijau.
		\item[PERENCANAAN KOTA] Tata ruang dan kebijakan pembangunan.
	\end{description}

	\vspace{0.5cm}
	\textbf{Keempatnya harus bekerja sebagai satu sistem tapak.}
\end{frame}

% ==========================================
% TOPIK 6: TUJUAN PERANCANGAN TAPAK
% ==========================================
\begin{frame}[mytitle=standard, t, mycolor=digiPH_white, light]
	\frametitle{Tujuan Perancangan Tapak}

	\begin{itemize}
		\item \textbf{Keterpaduan:} Keseluruhan program ruang dapat diwujudkan secara terpadu.
		\item \textbf{Kesesuaian:} Rancangan sesuai dengan potensi dan kendala tapak.
		\item \textbf{Kenyamanan:} Tapak menjadi wadah kegiatan yang nyaman.
		\item \textbf{Keamanan dan kesehatan:} Ruang mendukung keamanan serta kondisi lingkungan yang sehat.
		\item \textbf{Estetika:} Bangunan, ruang luar, dan lingkungan membentuk hubungan visual yang harmonis.
	\end{itemize}

	\vspace{0.3cm}
	\textit{Catatan: Tujuan perancangan dipengaruhi keterbatasan tapak dan dapat mencakup pertimbangan biaya maupun kapasitas.}
\end{frame}

% ==========================================
% TOPIK 7: KAIDAH PERANCANGAN TAPAK
% ==========================================
\begin{frame}[mytitle=center, c, mycolor=digiPH_leaf, dark]
	\frametitle{Kaidah Perancangan Tapak}
	\framesubtitle{Lima Kaidah Utama}

	\begin{itemize}
		\item \textbf{ORIENTASI:} Bagaimana bangunan dan elemen tapak diarahkan terhadap kondisi lingkungan.
		\item \textbf{AKSESIBILITAS:} Bagaimana pengguna masuk, bergerak, dan mencapai berbagai bagian tapak.
		\item \textbf{VEGETASI:} Bagaimana tumbuhan berperan secara ekologis, fungsional, dan visual.
		\item \textbf{VIEW:} Bagaimana kualitas pandangan dari dalam maupun ke luar tapak.
		\item \textbf{EKSISTING:} Bagaimana kondisi yang sudah ada dipertahankan, dimanfaatkan, atau direspons.
	\end{itemize}
\end{frame}

% ==========================================
% TOPIK 8: KAIDAH 1 (ORIENTASI)
% ==========================================
\begin{frame}[mytitle=standard, t, mycolor=digiPH_cream, light]
	\frametitle{Kaidah 1: Orientasi}

	Orientasi berkaitan erat dengan perwajahan (fasade) bangunan, perencanaan bukaan, tata letak denah, peletakan elemen arsitektur, dan hubungan bangunan terhadap matahari.

	\vspace{0.3cm}
	\textbf{Pertanyaan analisis:}
	\begin{itemize}
		\item Dari arah mana matahari datang?
		\item Bagaimana posisi bangunan terhadap matahari?
		\item Bagaimana bukaan ditempatkan?
		\item Bagaimana orientasi memengaruhi kenyamanan ruang?
	\end{itemize}

	\vspace{0.3cm}
	\textbf{Prinsip:} Orientasi bangunan tidak hanya menentukan arah hadap, tetapi juga memengaruhi kualitas ruang dalam dan ruang luar.
\end{frame}

% ==========================================
% TOPIK 9: KAIDAH 2 (AKSESIBILITAS)
% ==========================================
\begin{frame}[mytitle=standard, t, mycolor=digiPH_gray, light]
	\frametitle{Kaidah 2: Aksesibilitas}

	Perencanaan aksesibilitas memengaruhi alur atau \textit{flow} pengguna di dalam tapak.

	\vspace{0.3cm}
	\textbf{Hal yang perlu diperhatikan:}
	\begin{itemize}
		\item \textit{Main entrance} dan \textit{side entrance}
		\item Jalur kendaraan dan jalur pejalan kaki
		\item Hubungan antarruang dan akses visual
		\item Pemisahan area publik dan privat
	\end{itemize}

	\vspace{0.3cm}
	\textbf{Konsep penting:} Akses publik tidak seharusnya secara mudah mengintervensi area privat. Aksesibilitas adalah persoalan alur, hierarki, keamanan, dan privasi.
\end{frame}

% ==========================================
% TOPIK 10: KAIDAH 3 (VEGETASI)
% ==========================================
\begin{frame}[mytitle=standard, t, mycolor=digiPH_white, light]
	\frametitle{Kaidah 3: Vegetasi}

	\textbf{Fungsi Ekologis \& Teknis}
	\begin{itemize}
		\item Membentuk area teduh dan membantu konservasi energi.
		\item Mengurangi kebisingan dan membantu mencegah erosi.
	\end{itemize}

	\vspace{0.3cm}
	\textbf{Fungsi Visual \& Arsitektural}
	\begin{itemize}
		\item Menyaring pandangan dan memperindah pemandangan.
		\item Memperkuat hubungan bangunan dengan tapak.
	\end{itemize}

	\vspace{0.3cm}
	\textbf{Pertanyaan desain:} ``Apakah vegetasi hanya ditanam untuk mempercantik?'' \\
	\textit{Jawabannya: vegetasi merupakan elemen perancangan tapak yang bersifat fungsional sekaligus estetis.}
\end{frame}

% ==========================================
% TOPIK 11: KAIDAH 4 (VIEW)
% ==========================================
\begin{frame}[mytitle=standard, c, mycolor=digiPH_darkblue, dark]
	\frametitle{Kaidah 4: View}

	View berkaitan dengan kualitas kenyamanan visual penghuni. Perencanaan view perlu mempertimbangkan:
	\begin{itemize}
		\item pemandangan yang tersedia,
		\item penerangan alami dan pemanasan akibat radiasi matahari,
		\item arah bukaan serta privasi penghuni dan tetangga.
	\end{itemize}

	Ketika view luar tidak mendukung, pemandangan dibentuk melalui pengolahan taman, vegetasi, dan posisi bukaan.

	\vspace{0.3cm}
	\textbf{Prinsip:} View yang baik bukan sekadar indah, tetapi sesuai dengan kebutuhan fungsi dan privasi.
\end{frame}

% ==========================================
% TOPIK 12: KAIDAH 5 (EKSISTING)
% ==========================================
\begin{frame}[mytitle=standard, t, mycolor=digiPH_white, light]
	\frametitle{Kaidah 5: Eksisting}

	Kondisi eksisting harus diidentifikasi dan dimanfaatkan (contoh: rumah tetangga, jalan, pepohonan, pemandangan, kontur).

	\vspace{0.3cm}
	\textbf{Cara berpikir perancang:} \\
	Pertahankan $\rightarrow$ Manfaatkan $\rightarrow$ Adaptasikan $\rightarrow$ Sesuaikan

	\vspace{0.3cm}
	\textbf{Pertanyaan analisis:}
	\begin{itemize}
		\item Apa yang harus dipertahankan?
		\item Apa yang dapat dimanfaatkan?
		\item Apa yang menjadi kendala?
		\item Apa yang boleh diubah?
	\end{itemize}

	\textit{Eksisting bukan selalu hambatan; ia dapat menjadi aset perancangan.}
\end{frame}

% ==========================================
% TOPIK 13A: UNSUR-UNSUR PERENCANAAN (Sub 1 & 2)
% ==========================================
\begin{frame}[mytitle=standard, t, mycolor=digiPH_gray, light]
	\frametitle{Unsur-Unsur Perencanaan Tapak}
	\framesubtitle{Bagian 1: Akses dan Organisasi Ruang}

	Unsur tapak sangat luas dan saling berhubungan.

	\vspace{0.3cm}
	\textbf{Akses dan Sirkulasi}
	\begin{itemize}
		\item \textit{Main entrance} dan \textit{side entrance}
		\item Sirkulasi dan area parkir
	\end{itemize}

	\vspace{0.3cm}
	\textbf{Organisasi Ruang}
	\begin{itemize}
		\item Tata guna lahan dan \textit{zoning}
		\item Pola ruang luar
		\item Bukaan dan tutupan ruang luar
	\end{itemize}
\end{frame}

% ==========================================
% TOPIK 13B: UNSUR-UNSUR PERENCANAAN (Sub 3 & 4)
% ==========================================
\begin{frame}[mytitle=standard, t, mycolor=digiPH_cream, light]
	\frametitle{Unsur-Unsur Perencanaan Tapak}
	\framesubtitle{Bagian 2: Massa Bangunan dan Utilitas}

	\textbf{Massa dan Bangunan}
	\begin{itemize}
		\item Bentuk massa dan bentuk bangunan
		\item Orientasi dan komposisi massa
		\item Bukaan dan tutupan massa
	\end{itemize}

	\vspace{0.3cm}
	\textbf{Kenyamanan dan Utilitas}
	\begin{itemize}
		\item Penghawaan dan penerangan
		\item Jaringan komunikasi
	\end{itemize}

	\vspace{0.3cm}
	\textit{Semua unsur tersebut perlu dipertimbangkan sebagai satu kesatuan, bukan sebagai komponen yang terpisah.}
\end{frame}

% ==========================================
% TOPIK 14: FAKTOR ALAM (ANGIN)
% ==========================================
\begin{frame}[mytitle=standard, c, mycolor=digiPH_ocean, dark]
	\frametitle{Faktor Alam: Angin}

	\begin{description}
		\item[Angin laminer] Aliran angin berlapis-lapis dengan arah dan kecepatan yang relatif tidak berubah antar lapisan.
		\item[Angin terpisah] Terjadi ketika perubahan topografi menyebabkan perubahan pola aliran dan pemisahan lapisan angin.
		\item[Angin turbulen] Perubahan pola aliran menghasilkan turbulensi, terutama akibat perubahan topografi (bukit dan lembah).
	\end{description}

	\vspace{0.5cm}
	\textbf{Implikasi dalam desain:} \\
	Kontur + arah angin + posisi bangunan = memengaruhi kualitas aliran udara pada tapak.
\end{frame}

% ==========================================
% TOPIK 15A: FAKTOR ALAM (TANAH)
% ==========================================
\begin{frame}[mytitle=standard, t, mycolor=digiPH_white, light]
	\frametitle{Faktor Alam: Tanah dan Topografi}
	\framesubtitle{A. Analisis Tanah}

	Hal-hal yang perlu dianalisis terkait kondisi tanah pada tapak:
	\begin{itemize}
		\item Jenis dan kondisi tanah.
		\item Perubahan jenis tanah dalam tapak.
		\item Tingkat keasaman atau kebasaan.
		\item Ketebalan lapisan humus.
		\item Kemampuan tanah dalam menyerap air.
		\item Kemampuan tanah dalam mencegah erosi.
	\end{itemize}
\end{frame}

% ==========================================
% TOPIK 15B: FAKTOR ALAM (TOPOGRAFI)
% ==========================================
\begin{frame}[mytitle=standard, t, mycolor=digiPH_gray, light]
	\frametitle{Faktor Alam: Tanah dan Topografi}
	\framesubtitle{B. Analisis Topografi}

	Hal-hal yang perlu diperhatikan terkait topografi tapak:
	\begin{itemize}
		\item Kecuraman atau kedataran lereng.
		\item Keseragaman permukaan tanah.
		\item Hubungan dengan permukaan di sekitar tapak.
		\item Elemen alam yang tidak dapat diubah.
		\item Risiko terjadinya erosi.
		\item Orientasi lereng.
	\end{itemize}

	\vspace{0.3cm}
	\textbf{Kesimpulan:} Tanah dan topografi memengaruhi letak bangunan, sirkulasi, ruang luar, dan pengolahan permukaan tapak.
\end{frame}

% ==========================================
% TOPIK 16: KONTUR & CUT AND FILL
% ==========================================
\begin{frame}[mytitle=standard, c, mycolor=digiPH_leaf, dark]
	\frametitle{Kontur dan Konsep Cut and Fill}

	Kontur menggambarkan kondisi dan relief permukaan tapak.
	\begin{itemize}
		\item \textbf{CUT (KUPASAN):} Proses penggalian atau pemotongan tanah.
		\item \textbf{FILL (URUGAN):} Proses penimbunan atau penambahan tanah.
		\item \textbf{CUT AND FILL:} Kombinasi pemotongan dan penambahan tanah untuk memperoleh ketinggian permukaan yang sesuai.
	\end{itemize}

	\vspace{0.3cm}
	\textbf{Pertanyaan desain:} ``Apakah seluruh tapak harus dibuat datar?'' \\
	Tidak selalu. Tujuannya bukan sekadar meratakan tanah, tetapi menghasilkan kondisi tapak yang sesuai fungsi dan lingkungannya.
\end{frame}

% ==========================================
% TOPIK 17A: INTEGRASI (ALUR BERPIKIR)
% ==========================================
\begin{frame}[mytitle=center, t, mycolor=digiPH_white, light]
	\frametitle{Integrasi: Dari Analisis Menuju Perancangan}
	\framesubtitle{Alur Berpikir Perancangan Tapak}

	\begin{itemize}
		\item \textbf{BACA TAPAK:} Identifikasi kondisi alam, fisik, dan sosial.
		\item \textbf{TEMUKAN POTENSI \& KENDALA:} Apa yang dapat dimanfaatkan? Apa yang harus dihindari?
		\item \textbf{ANALISIS KAIDAH TAPAK:} Orientasi — Aksesibilitas — Vegetasi — View — Eksisting.
		\item \textbf{ANALISIS FAKTOR ALAM:} Angin — Tanah — Topografi — Kontur.
		\item \textbf{SUSUN PROGRAM \& ZONING:} Publik — Semipublik — Privat — Servis.
		\item \textbf{INTEGRASIKAN:} Bangunan + sirkulasi + ruang luar + utilitas + lingkungan sekitar.
	\end{itemize}
\end{frame}

% ==========================================
% TOPIK 17B: INTEGRASI (HASIL & REFLEKSI)
% ==========================================
\begin{frame}[mytitle=standard, c, mycolor=digiPH_cream, light]
	\frametitle{Integrasi: Dari Analisis Menuju Perancangan}
	\framesubtitle{Hasil Akhir dan Refleksi}

	\textbf{HASIL AKHIR:} \\
	Tapak yang fungsional, nyaman, aman, sehat, estetis, dan terpadu.

	\vspace{0.8cm}
	\begin{alertblock}{Refleksi Akhir}
		\centering
		\textit{“Perancangan tapak yang baik bukan dimulai dari menggambar bangunan, tetapi dari memahami tempat.”}
	\end{alertblock}
\end{frame}

% ==========================================
% TOPIK 17C: INTEGRASI (LATIHAN)
% ==========================================
\begin{frame}[mytitle=standard, c, mycolor=digiPH_darkblue, dark]
	\frametitle{Integrasi: Dari Analisis Menuju Perancangan}
	\framesubtitle{Latihan Singkat}

	Pilih satu tapak di sekitar Anda. Identifikasi:
	\begin{itemize}
		\item satu potensi dan satu kendala,
		\item arah orientasi dan akses utama,
		\item vegetasi yang dapat dipertahankan,
		\item potensi view dan kondisi kontur,
		\item serta faktor angin yang mungkin memengaruhi tapak.
	\end{itemize}

	\vspace{0.5cm}
	\scriptsize{\textit{Materi sumber menempatkan identifikasi potensi-kendala, kaidah perancangan, unsur tapak, serta faktor alam sebagai bagian penting dalam memahami perancangan tapak.}}
\end{frame}


\end{document}
